\documentclass[a4paper,12pt,twoside]{memoir}

% Castellano
\usepackage[spanish,es-tabla]{babel}
\selectlanguage{spanish}
\usepackage[utf8]{inputenc}
\usepackage[T1]{fontenc}
\usepackage{lmodern} % scalable font
\usepackage{microtype}
\usepackage{placeins}

%\usepackage{longtable}
\RequirePackage{booktabs}
\RequirePackage[table]{xcolor}
\RequirePackage{xtab}
\RequirePackage{multirow}

% Links
\PassOptionsToPackage{hyphens}{url}\usepackage[colorlinks]{hyperref}
\hypersetup{
	allcolors = {red}
}

% Más subseción
% Definición de \subsubsubsection
\newcommand{\subsubsubsection}[1]{
  \noindent\textbf{\quad #1}% Imprime el número de la sección y el título en negrita
  \par
}


% Ecuaciones
\usepackage{amsmath}

% Rutas de fichero / paquete
\newcommand{\ruta}[1]{{\sffamily #1}}

% Párrafos
\nonzeroparskip

% Huérfanas y viudas
\widowpenalty100000
\clubpenalty100000

% Evitar solapes en el header
\nouppercaseheads

% Imagenes
\usepackage{graphicx}
\newcommand{\imagen}[2]{
	\begin{figure}[!h]
		\centering
		\includegraphics[width=0.75\textwidth]{#1}
		\caption{#2}\label{fig:#1}
	\end{figure}
	\FloatBarrier
}

\newcommand{\imagenflotante}[2]{
	\begin{figure}%[!h]
		\centering
		\includegraphics[width=0.8\textwidth]{#1}
		\caption{#2}\label{fig:#1}
	\end{figure}
}



% El comando \figura nos permite insertar figuras comodamente, y utilizando
% siempre el mismo formato. Los parametros son:
% 1 -> Porcentaje del ancho de página que ocupará la figura (de 0 a 1)
% 2 --> Fichero de la imagen
% 3 --> Texto a pie de imagen
% 4 --> Etiqueta (label) para referencias
% 5 --> Opciones que queramos pasarle al \includegraphics
% 6 --> Opciones de posicionamiento a pasarle a \begin{figure}
\newcommand{\figuraConPosicion}[6]{%
  \setlength{\anchoFloat}{#1\textwidth}%
  \addtolength{\anchoFloat}{-4\fboxsep}%
  \setlength{\anchoFigura}{\anchoFloat}%
  \begin{figure}[#6]
    \begin{center}%
      \Ovalbox{%
        \begin{minipage}{\anchoFloat}%
          \begin{center}%
            \includegraphics[width=\anchoFigura,#5]{#2}%
            \caption{#3}%
            \label{#4}%
          \end{center}%
        \end{minipage}
      }%
    \end{center}%
  \end{figure}%
}

%
% Comando para incluir imágenes en formato apaisado (sin marco).
\newcommand{\figuraApaisadaSinMarco}[5]{%
  \begin{figure}%
    \begin{center}%
    \includegraphics[angle=90,height=#1\textheight,#5]{#2}%
    \caption{#3}%
    \label{#4}%
    \end{center}%
  \end{figure}%
}
% Para las tablas
\newcommand{\otoprule}{\midrule [\heavyrulewidth]}
%
% Nuevo comando para tablas pequeñas (menos de una página).
\newcommand{\tablaSmall}[5]{%
 \begin{table}
  \begin{center}
   \rowcolors {2}{gray!35}{}
   \begin{tabular}{#2}
    \toprule
    #4
    \otoprule
    #5
    \bottomrule
   \end{tabular}
   \caption{#1}
   \label{tabla:#3}
  \end{center}
 \end{table}
}

%
%Para el float H de tablaSmallSinColores
\usepackage{float}

%
% Nuevo comando para tablas pequeñas (menos de una página).
\newcommand{\tablaSmallSinColores}[5]{%
 \begin{table}[H]
  \begin{center}
   \begin{tabular}{#2}
    \toprule
    #4
    \otoprule
    #5
    \bottomrule
   \end{tabular}
   \caption{#1}
   \label{tabla:#3}
  \end{center}
 \end{table}
}

\newcommand{\tablaApaisadaSmall}[5]{%
\begin{landscape}
  \begin{table}
   \begin{center}
    \rowcolors {2}{gray!35}{}
    \begin{tabular}{#2}
     \toprule
     #4
     \otoprule
     #5
     \bottomrule
    \end{tabular}
    \caption{#1}
    \label{tabla:#3}
   \end{center}
  \end{table}
\end{landscape}
}

%
% Nuevo comando para tablas grandes con cabecera y filas alternas coloreadas en gris.
\newcommand{\tabla}[6]{%
  \begin{center}
    \tablefirsthead{
      \toprule
      #5
      \otoprule
    }
    \tablehead{
      \multicolumn{#3}{l}{\small\sl continúa desde la página anterior}\\
      \toprule
      #5
      \otoprule
    }
    \tabletail{
      \hline
      \multicolumn{#3}{r}{\small\sl continúa en la página siguiente}\\
    }
    \tablelasttail{
      \hline
    }
    \bottomcaption{#1}
    \rowcolors {2}{gray!35}{}
    \begin{xtabular}{#2}
      #6
      \bottomrule
    \end{xtabular}
    \label{tabla:#4}
  \end{center}
}

%
% Nuevo comando para tablas grandes con cabecera.
\newcommand{\tablaSinColores}[6]{%
  \begin{center}
    \tablefirsthead{
      \toprule
      #5
      \otoprule
    }
    \tablehead{
      \multicolumn{#3}{l}{\small\sl continúa desde la página anterior}\\
      \toprule
      #5
      \otoprule
    }
    \tabletail{
      \hline
      \multicolumn{#3}{r}{\small\sl continúa en la página siguiente}\\
    }
    \tablelasttail{
      \hline
    }
    \bottomcaption{#1}
    \begin{xtabular}{#2}
      #6
      \bottomrule
    \end{xtabular}
    \label{tabla:#4}
  \end{center}
}

%
% Nuevo comando para tablas grandes sin cabecera.
\newcommand{\tablaSinCabecera}[5]{%
  \begin{center}
    \tablefirsthead{
      \toprule
    }
    \tablehead{
      \multicolumn{#3}{l}{\small\sl continúa desde la página anterior}\\
      \hline
    }
    \tabletail{
      \hline
      \multicolumn{#3}{r}{\small\sl continúa en la página siguiente}\\
    }
    \tablelasttail{
      \hline
    }
    \bottomcaption{#1}
  \begin{xtabular}{#2}
    #5
   \bottomrule
  \end{xtabular}
  \label{tabla:#4}
  \end{center}
}



\definecolor{cgoLight}{HTML}{EEEEEE}
\definecolor{cgoExtralight}{HTML}{FFFFFF}

%
% Nuevo comando para tablas grandes sin cabecera.
\newcommand{\tablaSinCabeceraConBandas}[5]{%
  \begin{center}
    \tablefirsthead{
      \toprule
    }
    \tablehead{
      \multicolumn{#3}{l}{\small\sl continúa desde la página anterior}\\
      \hline
    }
    \tabletail{
      \hline
      \multicolumn{#3}{r}{\small\sl continúa en la página siguiente}\\
    }
    \tablelasttail{
      \hline
    }
    \bottomcaption{#1}
    \rowcolors[]{1}{cgoExtralight}{cgoLight}

  \begin{xtabular}{#2}
    #5
   \bottomrule
  \end{xtabular}
  \label{tabla:#4}
  \end{center}
}




\graphicspath{ {./img/} }

% Capítulos
\chapterstyle{bianchi}
\newcommand{\capitulo}[2]{
	\setcounter{chapter}{#1}
	\setcounter{section}{0}
	\setcounter{figure}{0}
	\setcounter{table}{0}
	\chapter*{#2}
	\addcontentsline{toc}{chapter}{#2}
	\markboth{#2}{#2}
}

% Apéndices
\renewcommand{\appendixname}{Apéndice}
\renewcommand*\cftappendixname{\appendixname}

\newcommand{\apendice}[1]{
	%\renewcommand{\thechapter}{A}
	\chapter{#1}
}

\renewcommand*\cftappendixname{\appendixname\ }

% Formato de portada
\makeatletter
\usepackage{xcolor}
\newcommand{\tutor}[1]{\def\@tutor{#1}}
\newcommand{\course}[1]{\def\@course{#1}}
\definecolor{cpardoBox}{HTML}{E6E6FF}
\def\maketitle{
  \null
  \thispagestyle{empty}
  % Cabecera ----------------
\noindent\includegraphics[width=\textwidth]{cabecera}\vspace{1cm}%
  \vfill
  % Título proyecto y escudo informática ----------------
  \colorbox{cpardoBox}{%
    \begin{minipage}{.8\textwidth}
      \vspace{.5cm}\Large
      \begin{center}
      \textbf{TFG del Grado en Ingeniería Informática}\vspace{.6cm}\\
      \textbf{\LARGE\@title{}}
      \end{center}
      \vspace{.2cm}
    \end{minipage}

  }%
  \hfill\begin{minipage}{.20\textwidth}
    \includegraphics[width=\textwidth]{escudoInfor}
  \end{minipage}
  \vfill
  % Datos de alumno, curso y tutores ------------------
  \begin{center}%
  {%
    \noindent\LARGE
    Presentado por \@author{}\\ 
    en Universidad de Burgos --- \@date{}\\
    Tutor: \@tutor{}\\
  }%
  \end{center}%
  \null
  \cleardoublepage
  }
\makeatother


% Datos de portada
\title{título del TFG \\Documentación Técnica}
\author{Daniel Fernández Barrientos}
\tutor{Raúl Marticorena Sanchez}
\date{\today}

\begin{document}

\maketitle



\cleardoublepage



%%%%%%%%%%%%%%%%%%%%%%%%%%%%%%%%%%%%%%%%%%%%%%%%%%%%%%%%%%%%%%%%%%%%%%%%%%%%%%%%%%%%%%%%



\frontmatter


\clearpage

% Indices
\tableofcontents

\clearpage

\listoffigures

\clearpage

\listoftables

\clearpage

\mainmatter

\appendix

\apendice{Plan de Proyecto Software}

\section{Introducción}

La fase de planificación es parte fundamental y necesaria en todos los proyectos.

Esta fase sive para ir evaluando la viabilidad del proyecto a lo largo de su vida, pudiendo prevenir o adelantarse a posibles inconvenientes que surgan durante la fase de desarrollo.
Tener medidas de forma correcta las tareas que se van realizando y su implicación en el total del proyecto es fundamental para un desarrollo sostenible, tanto económica como temporalmente.

Se ha decidido dividir esta fase en dos etapas:

\begin{itemize}
	\item
	Planificación temporal.
	\item
	Estudio de viabilidad.
\end{itemize}

La planificación temporal tratará sobre la gestión y aprovechamiento del tiempo, así como la división del proyecto en pequeños hitos, siguiendo una metodología ágil.

El estudio de viabilidad se descompone a su vez en:

\begin{itemize}
	\item
	Viabilidad económica: Inversión necesaria para acometer el proyecto y posible obtención de beneficios.
	\item
	Viabilidad legal: Todos los programas software conyevan unas implicaciones legales.
\end{itemize}

\section{Planificación temporal}

Para llevar a cabo una correcta planificación temporal se gestionado el proyecto siguiendo la metodología ágil que plantea Scrum, (poner enlace y bibliografía a Scrum) de forma "reducida" ya que el equipo de desarrolladores solo ha tenido un integrante, y las labores de "Scrum Master" y "Product Owner" han recaído en el tutor del TFG, el Sr. Raúl Marticorena.

Para el seguimiento del desarrollo incremental del proyecto software se ha divido la planificación en una serie de Sprints (enlace y referencia).

Para la visualización de los sprint se tenía intención de incluir los gráficos de "Burndown" pero debido a que han sido unos sprint vivos, donde se han inco incluyendo o desglosando tareas más genéricas en partes más específicas, se ha decidido utilizar los gráficos "Burnup" que representan de forma más real el número total de tareas completadas en cada sprint.

\subsection{Reunión de Kick-off}
\textit{Del 12/02/2024 - 1 hora}

La reunión de kick-off fue la toma de contacto con el proyecto, donde se trataron los aspectos básicos y se definieron varias herramientas para el desarrollo del mismo.

Temas tratados durante el Kick-off:
\begin{itemize}
	\item
	Análisis iniciald e la página de risk.real.eu.
	\item
	Borradores con la estrucutra de los cuestionarios.
	\item
	Valorar la carga de los cuestionarios en bases de datos o en ficheros de tipo json.
	\item
	Definición básica del proyecto:
	\begin{enumerate}
		\item
		Web que lea los cuestionarios.
		\item
		Usuarios separados por sesiones.
		\item
		Creación de cuestionarios desde la web.
	\end{enumerate}
	\item
	Diseño responsive: Permite el renderizado para otros dispositivos.
	\item
	Documentación y gestión adecuada del proyecto.
	\item
	Uso de Spring Boot.
	\item
	Github o ZenHub para la gestión de tareas.
\end{itemize}


\subsection{Sprint 1 - Base del proyecto}
\textit{Del 27/02/2024 al 11/03/2024 - 20 horas}

Se estima una dedicación total de 20 horas de trabajo.
Se cumple con la estimación inicial de horas para completar todas las tareas.

Temas tratados durante el Sprint:

\begin{itemize}
	\item
	Uso e integración de zube.io con GitHub.
	\item
	Registro en https://app.riskreal.eu/ para analizar la parte "privada" de la aplicación.
	\item
	Schema json:
	\begin{itemize}
		\item
		Posibilidad de imagen en cada respuesta
		\item
		Atributo idioma en el cuestionario
	\end{itemize}
	Empezar a valorar y tener en cuenta la posibilidad de cuestionarios "Multi-idioma".
	\item
	Investigar sobre Thumeleaf y Spring Tools.
	\item
	Cargar un json de prueba en la app.
	\item
	Representación en una web del json.
	\item
	Posible estructura de directorios de la aplicación.
	\item
	Abordar la generación dinámica del contenido respecto al json.
\end{itemize}

Tareas tratadas cerradas a lo largo del sprint:
\imagen{Sprint1_issues}{Sprint 1 - Base del proyecto - Tareas}

Gráfico de Burnup del sprint:
\imagen{Sprint1_burnup}{Sprint 1 - Base del proyecto - Gráfico}


\subsection{Sprint 2 - Evolución del prototipo}
\textit{Del 11/03/2024 al 25/03/2024 - 20 horas}

Se estima una dedicación total de 20 horas de trabajo.
Se cumple con la estimación inicial de horas para completar todas las tareas.

Temas tratados durante el Sprint:

\begin{itemize}
	\item
	Inclusión de más preguntas y respuestas en el cuestionario.
	\item
	Tratar las preguntas como independientes en la web (permitir transición entre preguntas).
	\item
	Recoger y mostrar valoraciones al terminar el cuestionario.
	\item
	Concepto de sesión web, para que no se mezclen datos de 2 usuarios.
	\item
	Primer acercamiento a un framework de estilo, bootstrap o similar.
\end{itemize}

Tareas tratadas cerradas a lo largo del sprint:
\imagen{Sprint2_issues}{Sprint 2 - Evolución del prototipo - Tareas}

Gráfico de Burnup del sprint:
\imagen{Sprint2_burnup}{Sprint 2 - Evolución del prototipo - Gráfico}


\subsection{Sprint 3 - Evolución del prototipo 2}
\textit{Del 25/03/2024 al 24/04/2024 - 60 horas}

Se estima una dedicación total de 20 horas de trabajo.
Tras encontrar problemas diversos problemas a la hora de integrar la seguridad de Spring Security "referencia" en la app, no solo no se cumple con la planificación inicial sino que se extiende el mismo spring dos semanas más, ampliando la carga de trabajo con más tareas y teniendo que emplear finalmente 60 horas de trabajo para cumplir con casi todas las tareas.
La tarea que queda pendiente se solucionará en el sprint 5.

Temas tratados durante el Sprint:
\begin{itemize}
	\item
	Independizar las sesiones de los usuarios.
	\item
	Integración de la seguridad en la aplicación.
	\item
	Integrar la funcionalidad de envío de correo electrónico desde la web.
	\item
	Enviar la información del cuestionario con el método POST para que no vayan los datos en claro.
	\item
	Internacionalización de los distintos menús de la app web.
	\item
	Integración de framework CSS para utilizar un diseño "responsive" en la web.
	\item
	Recoger datos de puntuación de las preguntas.
	\item
	Mostrar un score tras finalizar el cuestionario.
\end{itemize}

Tareas tratadas cerradas a lo largo del sprint:
\imagen{Sprint3_issues}{Sprint 3 - Evolución del prototipo 2 - Tareas}

Gráfico de Burnup del sprint:
\imagen{Sprint3_burnup}{Sprint 3 - Evolución del prototipo 2 - Gráfico}


\subsection{Sprint 4 - Evolución y finalización de la parte del cuestionario}
\textit{Del 24/04/2024 al 06/05/2024 - 23 horas}

Se estima una duración inicial de 20 horas de trabajo, aunque finalmente se emplean 23 horas para cumplir con 4 de las 6 tareas previstas.
La tarea de implementar el registro de usuarios a través de una base de datos llevó más tiempo del esperado, de ahí la necesidad de horas adicionales y el no poder cumplir con todas las tareas previstas.
Las 2 tareas pendientes se completarán en el sprint 5.

Temas tratados durante el Sprint:
\begin{itemize}
	\item
	Gestión de usuarios y roles.
	\item
	Mostrar información del usuario logueado.
	\item
	Creación de un formulario de registro de usuarios.
	\item
	Almacenar los datos de score en un fichero, para futuros análisis.
	\item
	Introducción al uso de roles para acceder a distintas partes de la aplicación.
\end{itemize}

Tareas tratadas cerradas a lo largo del sprint:
\imagen{Sprint4_issues}{Sprint 4 - Evolución y finalización de la parte del cuestionario - Tareas}

Gráfico de Burnup del sprint:
\imagen{Sprint4_burnup}{Sprint 4 - Evolución y finalización de la parte del cuestionario - Gráfico}

\subsection{Sprint 5 - Evolución y finalización de la parte del cuestionario 2}
\textit{Del 06/05/2024 al 22/05/2024 - 40 horas}

Se estima una dedicación total de 40 horas de trabajo.
Se cumple con la estimación inicial de horas para completar todas las tareas.
En este sprint se han completado las tareas que quedaron pendientes del sprint 3 y 4, por eso la carga de trabajo y de horas ha sido mayor que en los anteriores.

Temas tratados durante el Sprint:

\begin{itemize}
	\item
	Investigar sitio para alojar un test de la aplicación --> https://www.heroku.com/.
	Guardar el score por pregunta y el total.
	\item
	Descarga de resultados en CSV.
	\item
	Array de imágenes en las preguntas.
	\item
	Diseño de datos bbdd y json. Código de la aplicación.
	\item
	Revisar información para incluir en el arquitectónico.
	\item
	Revisar diagramas para incluir en el diseño procedimental.
\end{itemize}

Tareas tratadas cerradas a lo largo del sprint:
\imagen{Sprint5_issues}{Sprint 5 - Evolución y finalización de la parte del cuestionario 2 - Tareas}

Gráfico de Burnup del sprint:
\imagen{Sprint5_burnup}{Sprint 5 - Base del proyecto - Gráfico}

\subsection{Sprint 6 - Carga de cuestionarios y documentación}
\textit{Del 22/05/2024 al 30/05/2024 - 25 horas}

Temas tratados durante el Sprint:
\begin{itemize}
	\item
	Carga de cuestionarios desde una carpeta definida en configuración.
	\item
	Carga de cuestionarios de forma "viva" desde la propia web de la aplicación.
	\item
	Completar la memoria del proyecto.
	\item
	Completar los anexos del proyecto.
	\item
	Documentar de forma adecuada y en inglés el código de java.
\end{itemize}

Tareas tratadas cerradas a lo largo del sprint:
% \imagen{Sprint6_issues}{Sprint 6 - Carga de cuestionarios y documentación - Tareas}

Gráfico de Burnup del sprint:
% \imagen{Sprint6_burnup}{Sprint 6 - Carga de cuestionarios y documentación - Gráfico}


\subsection{Sprint 7 - }
\textit{Del 30/05/2024 al 03/06/2024 - 20 horas}

Temas tratados durante el Sprint:
\begin{itemize}
	\item
	Lista de tareas.
\end{itemize}

Tareas tratadas cerradas a lo largo del sprint:
% \imagen{Sprint7_issues}{Sprint 7 -  - Tareas}

Gráfico de Burnup del sprint:
% \imagen{Sprint7_burnup}{Sprint 7 -  - Gráfico}


\subsection{Sprint 8 - }
\textit{Del 03/05/2024 al 10/06/2024 - 30 horas}

Temas tratados durante el Sprint:
\begin{itemize}
	\item
	Lista de tareas.
\end{itemize}

Tareas tratadas cerradas a lo largo del sprint:
% \imagen{Sprint8_issues}{Sprint 8 -  - Tareas}

Gráfico de Burnup del sprint:
% \imagen{Sprint8_burnup}{Sprint 8 -  - Gráfico}

\section{Estudio de viabilidad}

\subsection{Viabilidad económica}

A lo largo de este apartado se van a simular los costes/beneficios del proyecto, si hubiese sido llevado a cabo en una empresa real.

\subsubsection{Costes}
Los costes se han desglosado en dos categorías:
\begin{itemize}
	\item	Costes de personal
	\item	Costes de recursos
\end{itemize}

\subsubsubsection{Costes de personal}
Se ha simulado la contratación de un programador web durante los 5 meses que ha durado el proyecto.

En la siguiente tabla se desglosan los datos de retenciones "referencia" y tributación de la empresa a la seguridad social "referencia" para hacer la estimación.
%URL calculo nómina https://www.billin.net/calculadora-contratar-trabajador/

%URL IRPF https://cincodias.elpais.com/herramientas/calculadora-irpf/

%URL Cotización Régimen General seguridad social: https://www.seg-social.es/wps/portal/wss/internet/Trabajadores/CotizacionRecaudacionTrabajadores/10721/10957/583#580.

\tablaSmallSinColores{Coste de personal}{l c c}{costepersonal}
{ \multicolumn{1}{l}{\textbf{Concepto}} & \textbf{Coste}\\}{ 
Salario mensual bruto & 2.000,00 € \\
Retención IRPF (13,55\%) & 217,00 € \\
Seguridad social (36,25\%) & 1.450,00 € \\
Salario mensual neto & 1.602,00 € \\
Coste total para la empresa & 3.323,00 €\\
\midrule
\textbf{Total 5 meses}  & \textbf{16.615,00 €}\\
} 
\subsubsubsection{Costes de recursos}
Solo se ha tenido en cuenta el material necesario para el nuevo empleado, no se considera ni el alquiler del espacio, ni el acceso a internet ni el mobiliario, ya que gracias al trabajo híbrido (presencial + teletrabajo) todo esto no ha supuesto ningún gasto adicional para la empresa.

A continuación se detalla de forma breve el material hardware que se ha tenido en cuenta:
\begin{itemize}
	% https://www.dell.com/es-es/shop/laptops/nuevo-inspiron-14-plus/spd/inspiron-14-7440-laptop/cn74702
	\item Ordenador portátil: Dell Inspiron 14 Plus 1.098,99 € - 908,26 € sin IVA
	%https://www.pccomponentes.com/dell-s-series-s2721hs-27-led-ips-fullhd-freesync
	\item	Monitor adicional: Dell S Series S2721HS 27" LED IPS FullHD  229,36€  - 189,55 € SIN IVA
	%https://www.pccomponentes.com/logitech-mk235-combo-teclado-y-raton-inalambrico
	\item	Combo teclado y ratón: Logitech MK235 39,99 € - 33,05€ SIN IVA
	\item	Licencia Windows 11 pro: Incluida en el precio del portátil
\end{itemize}

Debido a que las empresas se deducen el IVA, los costes 

\tablaSmallSinColores{Coste de recursos}{l c c}{costerecursos}
{ \multicolumn{1}{l}{\textbf{Concepto}} & \textbf{Coste}\\}{ 
Ordenador portátil & 908,26 € \\
Monitor adicional  & 189,55 € \\
Combo teclado y ratón & 33,05€ \\
S.O. Windows 11 Pro & - € \\
\midrule
\textbf{Total}  & \textbf{1.130,86 €}\\
} 


\tablaSmallSinColores{Coste Total}{l c c}{costetotall}
{ \multicolumn{1}{l}{\textbf{Concepto}} & \textbf{Coste}\\}{ 
Coste de personal & 16.615,00 € \\
Coste de material  & 1.130,86 € \\
\midrule
\textbf{Total}  & \textbf{17.745,86 €}\\
} 


\subsection{Viabilidad legal}



\apendice{Especificación de Requisitos}

\section{Introducción}

En este anexo se recoge la especificación de requisitos de la aplicación web que se ha desarrollado.

Este documento tiene el objetivo doble de servir de \textit{contrato} con el cliente y como base de documentación para el análisis de la aplicación.

\section{Objetivos generales}

El proyecto tiene los siguientes objetivos generales:

\begin{itemize}
	\item Desarrollar una aplicación web que permita la realización de cuestionarios para evaluar \textit{Soft Skills}.
	\item Permitir la carga de cuestionarios de forma dinámica, tanto a través de una carpeta como desde la propia interfaz web.
	\item Almacenar los resultados de los cuestionarios de forma anónima con fines estadísticos para los creadores.
	\item La descarga completa de los resultados en un documento correctamente estructurado y que sea sencillo de modelar.
	\item La aplicación web tiene que tener un diseño \textit{responsive} que extienda su funcionalidad a diferentes dispositivos.
	\item La aplicación tiene que tener ser multi-idioma, para poder llegar a más usuarios.
\end{itemize}

\section{Catálogo de requisitos}

Se van a enumerar los requisitos separándolos en dos categorías: requisitos funcionales y no funcionales \cite{web:requisitos}.

\subsection{Requisitos funcionales}
\begin{itemize}
	\item \textbf{RF-1 Gestión de cuestionarios}: La aplicación web tiene que ser capaz de gestionar cuestionarios.
	\begin{itemize}
		\item \textbf{RF-1.1 Carga de cuestionario}: Se tienen que poder cargar nuevos cuestionarios en la aplicación web.
		\item \textbf{RF-1.2 Generación de cuestionario}: El administrador debe poder generar cuestionarios a partir de los archivos json cargados.
		\item \textbf{RF-1.3 Realización de cuestionario}: Los usuarios deben poder realizar los cuestionarios disponibles.
	\end{itemize}
	\item \textbf{RF-2 Gestión de usuarios}: La aplicación web tiene que ser capaz de gestionar usuarios.
		\begin{itemize}
		\item \textbf{RF-2.1 Registro de usuario}: Los usuarios deben poder registrarse en la aplicación.
		\item \textbf{RF-2.2 Modificación de datos}: Los usuarios deben poder editar sus datos.
	\end{itemize}
	\item \textbf{RF-3 Gestión de resultados}: La apicación web tiene que permitir la gestión de los resultados obtenidos.
		\begin{itemize}
		\item \textbf{RF-3.1 Almacenamiento de resultados}: Los resultados se tienen que almacenar de forma persistente en la aplicación.
		\item \textbf{RF-3.2 Descarga de resultados}: Los resultatos se tienen que poder descargar de la aplicación.
	\end{itemize}
	\item \textbf{RF-4 Diseño internacional}: La aplicación web tiene que ser capaz de adaptarse a varios idiomas.
		\begin{itemize}
		\item \textbf{RF-4.1 Selección de idioma}: La aplicación debe permitir la selección del idioma.
		\item \textbf{RF-4.2 Idioma de los cuestionarios}: La aplicación debe permitir la carga de cuestionarios en varios idiomas.
	\end{itemize}
\end{itemize}

\subsection{Requisitos no funcionales}
\begin{itemize}
	\item \textbf{RNF-1 Usabilidad}: La aplicación debe ser intuitiva, con una curva de aprendizaje baja.
	\item \textbf{RNF-2 Rendimiento}: La aplicación web tiene que tener unos tiempo de carga y respuesta bajos.
	\item \textbf{RNF-3 Escalabilidad}: La aplicación web tiene que poder funcionar en cualquier navegador, y permitir añadir funcionalidades de forma sencilla.
	\item \textbf{RNF-4 Portabilidad}: La aplicación web tiene que adaptarse a diferentes dispositivos.
	\item \textbf{RNF-5 Internacionalizaicón}: La aplicación web tiene que estar preparada para soportar varios idiomas.
	\item \textbf{RNF-6 Seguridad}: La aplicación web debe gestionar los accesos y datos sensibles de forma adecuada.
\end{itemize}

\section{Especificación de requisitos}

Esta sección va a incluir el diagrama de casos de uso y la explicación de cada uno.

\begin{landscape}
\subsection{Diagrama de casos de uso}
\imagenAncho{diagramaCasosUso}{Diagrama de casos de uso.}{1.35}
\end{landscape}

\subsection{Actores}

Pueden interactuar con el sistema 3 actores diferentes:

\begin{itemize}
	\item Usuario invitado: Usuario que accede a la página de forma anómina, y puede probar los cuestionarios de ejemplo o registrarse.  
	\item Usuario registrado: Usuario que accede a la página e inicia sesión. Puede ver su perfíl y modificarlo o realizar los cuestionarios disponibles.
	\item Administrador: Usuario administrador que puede gestionar roles de los usuarios registrados, descargarse los resultados completos de todos los cuestionarios realizados, cargar y generar más cuestionarios, o realizar alguno de los cuestionarios disponibles.
\end{itemize}

\begin{table}[p]
	\centering
	\begin{tabularx}{\linewidth}{ p{0.21\columnwidth} p{0.71\columnwidth} }
		\toprule
		\textbf{CU-01}    & \textbf{Cargar Cuestionarios}\\
		\toprule
		\textbf{Versión}              & 1.0    \\
		\textbf{Autor}                & Daniel Fernández Barrientos \\
		\textbf{Requisitos asociados} & RF-1, RF-1.1 \\
		\textbf{Descripción}          & Permite al administrador cargar nuevos cuestionarios. \\
		\textbf{Precondición}         & La aplicacion se encuentra levantada, y ha iniciado sesión un usuario administrador. \\
		\textbf{Acciones}             &
		\begin{enumerate}
			\def\labelenumi{\arabic{enumi}.}
			\tightlist
			\item El usuario administrador inicia sesión.
			\item Accede al apartador de carga de cuestionarios.
			\item Selecciona un archivo json con una estructura válida y lo carga en la aplicación.
		\end{enumerate}\\
		\textbf{Postcondición}        & El nuevo archivo json se encuentra disponible para generar un nuevo cuestionario. \\
		\textbf{Excepciones}          & Error si el archivo seleccionado no es un json o si su estructura no es la esperada. \\
		\textbf{Importancia}          & Alta. \\
		\bottomrule
	\end{tabularx}
	\caption{CU-01 Cargar Cuestionarios.}
\end{table}

\begin{table}[p]
	\centering
	\begin{tabularx}{\linewidth}{ p{0.21\columnwidth} p{0.71\columnwidth} }
		\toprule
		\textbf{CU-02}    & \textbf{Generar Cuestionarios}\\
		\toprule
		\textbf{Versión}              & 1.0    \\
		\textbf{Autor}                & Daniel Fernández Barrientos \\
		\textbf{Requisitos asociados} & RF-1, RF-1.2 \\
		\textbf{Descripción}          & Permite al administrador generar nuevos cuestionarios. \\
		\textbf{Precondición}         & La aplicacion se encuentra levantada, ha iniciado sesión un usuario administrador y hay cuestionarios pendientes de generar. \\
		\textbf{Acciones}             &
		\begin{enumerate}
			\def\labelenumi{\arabic{enumi}.}
			\tightlist
			\item El usuario administrador inicia sesión.
			\item Accede al apartador de carga de cuestionarios.
			\item Selecciona uno de los archivos json cargados y genera un nuevo cuestionario.
		\end{enumerate}\\
		\textbf{Postcondición}        & Hay un nuevo cuestionario disponible en la aplicación.. \\
		\textbf{Excepciones}          & Error si el archivo json utilizado no tiene una estructura correcta. \\
		\textbf{Importancia}          & Alta. \\
		\bottomrule
	\end{tabularx}
	\caption{CU-02 Generar Cuestionarios.}
\end{table}

\begin{table}[p]
	\centering
	\begin{tabularx}{\linewidth}{ p{0.21\columnwidth} p{0.71\columnwidth} }
		\toprule
		\textbf{CU-03}    & \textbf{Realizar Cuestionario}\\
		\toprule
		\textbf{Versión}              & 1.0    \\
		\textbf{Autor}                & Daniel Fernández Barrientos \\
		\textbf{Requisitos asociados} & RF-1.3, RF-3, RF-3.1 \\
		\textbf{Descripción}          & Permite a cualquier usuario realizar cuestionarios. \\
		\textbf{Precondición}         & La aplicacion se encuentra levantada y hay cuestionarios disponibles en el idioma seleccionado. \\
		\textbf{Acciones}             &
		\begin{enumerate}
			\def\labelenumi{\arabic{enumi}.}
			\tightlist
			\item Un usuario accede a la página web.
			\item Selecciona su idioma o mantiene el por defecto.
			\item Accede a uno de los cuestionarios que se le presentan.
			\item Contesta a las preguntas y llega hasta el final, donde termina el cuestionario.
			\item Ve una página con el resultado obtenido.
		\end{enumerate}\\
		\textbf{Postcondición}        & Se ha incluido una línea nueva en el csv de resultados. \\
		\textbf{Excepciones}          & Alguna de las imágenes del cuestionario no está presente. \\
		\textbf{Importancia}          & Alta. \\
		\bottomrule
	\end{tabularx}
	\caption{CU-03 Realizar Cuestionario.}
\end{table}

\begin{table}[p]
	\centering
	\begin{tabularx}{\linewidth}{ p{0.21\columnwidth} p{0.71\columnwidth} }
		\toprule
		\textbf{CU-04}    & \textbf{Registro de Usuarios}\\
		\toprule
		\textbf{Versión}              & 1.0    \\
		\textbf{Autor}                & Daniel Fernández Barrientos \\
		\textbf{Requisitos asociados} & RF-2, RF-2.1 \\
		\textbf{Descripción}          & Permite a cualquier usuario registrarse en la aplicación web. \\
		\textbf{Precondición}         & La aplicacion se encuentra levantada junto con la base de datos. \\
		\textbf{Acciones}             &
		\begin{enumerate}
			\def\labelenumi{\arabic{enumi}.}
			\tightlist
			\item Un usuario accede a la página web.
			\item Selecciona el registro.
			\item Completa el formulario y se registra.
		\end{enumerate}\\
		\textbf{Postcondición}        & Se ha añadido un nuevo usuario a la aplicación. \\
		\textbf{Importancia}          & Alta. \\
		\bottomrule
	\end{tabularx}
	\caption{CU-04 Registro de Usuarios.}
\end{table}

\begin{table}[p]
	\centering
	\begin{tabularx}{\linewidth}{ p{0.21\columnwidth} p{0.71\columnwidth} }
		\toprule
		\textbf{CU-05}    & \textbf{Iniciar Sesión en la web}\\
		\toprule
		\textbf{Versión}              & 1.0    \\
		\textbf{Autor}                & Daniel Fernández Barrientos \\
		\textbf{Requisitos asociados} & RF-2 \\
		\textbf{Descripción}          & Los usuarios registrados pueden iniciar sesión. \\
		\textbf{Precondición}         & La aplicacion se encuentra levantada y el usuario está registrado. \\
		\textbf{Acciones}             &
		\begin{enumerate}
			\def\labelenumi{\arabic{enumi}.}
			\tightlist
			\item Un usuario accede a la página web.
			\item Selecciona la opción de iniciar sesión.
			\item Introduce sus datos de usuario.
			\item Si son correctos, accede a la parte privada de la aplicación.
			\item Si no lo son, se permiten más intentos o la posibilidad de recibir un correo con una contraseña nueva de un solo uso.
			\item Se obliga al cambio tras introducir la contraseña recibida por correo.
		\end{enumerate}\\
		\textbf{Postcondición}        & Puede acceder a todos los cuestionarios disponibles y ver o actualizar su perfíl. \\
		\textbf{Excepciones}          & Usuario o contraseña incorrectos. \\
		\textbf{Importancia}          & Alta. \\
		\bottomrule
	\end{tabularx}
	\caption{CU-05 Iniciar Sesión en la web.}
\end{table}

\begin{table}[p]
	\centering
	\begin{tabularx}{\linewidth}{ p{0.21\columnwidth} p{0.71\columnwidth} }
		\toprule
		\textbf{CU-6}    & \textbf{Modificar datos del perfil}\\
		\toprule
		\textbf{Versión}              & 1.0    \\
		\textbf{Autor}                & Daniel Fernández Barrientos \\
		\textbf{Requisitos asociados} & RF-2, RF-2.2 \\
		\textbf{Descripción}          & Permite a cualquier usuario identificado modificar su perfil. \\
		\textbf{Precondición}         & El usuario ha iniciado sesión. \\
		\textbf{Acciones}             &
		\begin{enumerate}
			\def\labelenumi{\arabic{enumi}.}
			\tightlist
			\item Un usuario inicia sesión en la página web.
			\item Accede a su perfil.
			\item Modifica los datos que sean erróneos (el nombre de usuario no permite edición).
			\item Pulsa en actualizar y se muestra el perfil actualizado.
		\end{enumerate}\\
		\textbf{Postcondición}        & Se han actualizado sus datos de perfil. \\
		\textbf{Importancia}          & Alta. \\
		\bottomrule
	\end{tabularx}
	\caption{CU-06 Modificar datos del perfil.}
\end{table}

\begin{table}[p]
	\centering
	\begin{tabularx}{\linewidth}{ p{0.21\columnwidth} p{0.71\columnwidth} }
		\toprule
		\textbf{CU-07}    & \textbf{Gestión de Roles}\\
		\toprule
		\textbf{Versión}              & 1.0    \\
		\textbf{Autor}                & Daniel Fernández Barrientos \\
		\textbf{Requisitos asociados} & RF-2, RF-2.2 \\
		\textbf{Descripción}          & Permite a un administrador actualizar el rol en la aplicación de otros usuarios. \\
		\textbf{Precondición}         & La aplicación está levantada y hay usuarios registrados en la base de datos. \\
		\textbf{Acciones}             &
		\begin{enumerate}
			\def\labelenumi{\arabic{enumi}.}
			\tightlist
			\item Un administrador inicia sesión en la web.
			\item Selecciona la opción de gestión de roles.
			\item Modifica el perfíl de seguridad de los usuarios que considere.
			\item Pulsa en actualizar y se muestran de nuevos los perfiles actualizados.
		\end{enumerate}\\
		\textbf{Postcondición}        & Se ha actualizado el perfil de seguridad de los usuarios correspondientes. \\
		\textbf{Importancia}          & Media. \\
		\bottomrule
	\end{tabularx}
	\caption{CU-07. Gestión de Roles.}
\end{table}

\begin{table}[p]
	\centering
	\begin{tabularx}{\linewidth}{ p{0.21\columnwidth} p{0.71\columnwidth} }
		\toprule
		\textbf{CU-08}    & \textbf{Descarga de Resultados}\\
		\toprule
		\textbf{Versión}              & 1.0    \\
		\textbf{Autor}                & Daniel Fernández Barrientos \\
		\textbf{Requisitos asociados} & RF-3, RF-3.2 \\
		\textbf{Descripción}          & Permite a un administrador descargarse todos los resultados de los cuestionarios que se han realizado. \\
		\textbf{Precondición}         & La aplicación se encuentra levantada y ha iniciado sesión un administrador. \\
		\textbf{Acciones}             &
		\begin{enumerate}
			\def\labelenumi{\arabic{enumi}.}
			\tightlist
			\item Un administrador accede a la web.
			\item Selecciona la opción de descarga de resultados.
			\item Se descarga un csv con todos los resultados disponibles.
		\end{enumerate}\\
		\textbf{Postcondición}        & Se ha descargado un csv en el equipo del usuario administrador. \\
		\textbf{Excepciones}          & Si no se ha realizado ningún cuestionario, el csv estará vacío. \\
		\textbf{Importancia}          & Alta. \\
		\bottomrule
	\end{tabularx}
	\caption{CU-08 Descarga de Resultados.}
\end{table}

\begin{table}[p]
	\centering
	\begin{tabularx}{\linewidth}{ p{0.21\columnwidth} p{0.71\columnwidth} }
		\toprule
		\textbf{CU-09}    & \textbf{Selección de Idioma}\\
		\toprule
		\textbf{Versión}              & 1.0    \\
		\textbf{Autor}                & Daniel Fernández Barrientos \\
		\textbf{Requisitos asociados} & RF-4, RF-4.1. RF-4.2 \\
		\textbf{Descripción}          & Permite realizar cuestionarios o ver la aplicación en diferentes idiomas. \\
		\textbf{Precondición}         & La aplicación se encuentra levantada y hay cuestionarios generados en varios idiomas. \\
		\textbf{Acciones}             &
		\begin{enumerate}
			\def\labelenumi{\arabic{enumi}.}
			\tightlist
			\item Un usuario accede a la página web.
			\item Selecciona su idioma o mantiene el por defecto.
			\item Los menús de la web aparecen en el idioma seleccionado.
			\item Los cuestionarios aparecen en el idioma seleccionado.
		\end{enumerate}\\
		\textbf{Postcondición}        & La aplicación se muestra en el idioma seleccionado y le actualiza la lista de cuestionarios disponibles. \\
		\textbf{Excepciones}          & No hay cuestionarios diponibles en algún idioma. \\
		\textbf{Importancia}          & Alta. \\
		\bottomrule
	\end{tabularx}
	\caption{CU-09 Selección de Idioma.}
\end{table}



\apendice{Especificación de diseño}

\section{Introducción}

En este anexo se define la resolución de las especificaciones y objetivos expuestos anteriormente.

\section{Diseño de datos}

\subsection{Entidades}
La aplicación cuenta con las siguientes entidades de base de datos:

\begin{itemize}
	\item \textit{Users}: Tiene su correo electrónico, nombre, apellido, rango de edad, código de empresa, género, rol en la empresa, contraseña y el estado en la aplicación.
	\item \textit{Profiles}: Tiene el correo electrónico y la autoridad que tiene cada usuario en la aplicación.
\end{itemize}
\subsubsection{Diagrama Relacional}

La relación es 1 a 1 porque cada email solo puede estar una vez en cada tabla.

La tabla de \textit{Profiles} contiene la autoridad (permisos sobre la aplicación) de cada usuario dado de alta.

Actualmente solo se puede tener uno de estos dos permisos:

\begin{enumerate}
	\item \textit{ADMIN}: Usuario administrador con acceso a todas las opciones de la aplicación.
	 \item \textit{USER}: Usuario estándar de la aplicación, con acceso a cuestionarios y su perfil.
\end{enumerate}

\imagenB{diagramaClases}{Diagrama relacional}{.5}

\subsection{Clases}
Además de las entidades anteriores, se tienen las siguientes clases modelo, que contienen los objetos con los que trabaja la aplicación:

\begin{itemize}
	\item \textit{User}: Relacionada con la entidad \textit{Users}.
	\item \textit{Profile}: Relacionada con la entidad \textit{Profiles}.
	\item \textit{Quiz}: Tiene el id, valor máximo, título, idioma, descripción, tipo e imágenes propias del cuestionario y un listado de preguntas.
	\item \textit{Question}: Tiene el id, descripción, habilidad, e imágenes de la pregunta y un listado de respuestas.
	\item \textit{Answer}: Tiene el texto, la imagen, el id y el valor de cada respuesta.
	\item \textit{UserSelection}: Tiene los mapas de datos de respuestas, y de valores de las respuestas. Se utiliza para el cálculo del resultado tras la realización de un cuestionario.
	\item \textit{EmailRequest}: Tiene el campo del para, el asunto y el cuerpo del mensaje. Es la clase que se utiliza en el envío de correos electrónicos.
\end{itemize}
\subsubsection{Diagrama de Clases}
El siguiente diagrama es el diagrama de clases que sigue la estructura de los archivos \textit{JSON} de los cuestionarios que se cargan en la aplicación.
\imagenB{diagramaClasesJava}{Diagrama de clases}{.5}

\section{Diseño procedimental}
En este aspecto se detallan los aspectos del acceso a los cuestionarios.

La aplicación no realiza ningún proceso inicial de carga cuando arranca:

\begin{itemize}
	\item Los datos de los cuestionarios disponibles se muestran a la vez que se va accediendo a las diferentes páginas, para que sea transparente para lo usuarios la carga de nuevos cuestionarios.
	\item Cada vez que un usuario inicia sesión, se comprueba la base de datos.
\end{itemize}

Se muestra como ejemplo el diagrama de secuencia de un usuario que tras iniciar sesión, inicia un cuestionario.
\imagenB{diagramaAcceso}{Diagrama de secuencia en la aplicación web.}{.5}

\section{Diseño arquitectónico}

El hecho de que el proyecto se haya realizado utilizando \textit{Spring Boot} ha condicionado las decisiones del diseño.

\subsection{Modelo-Vista-Controlador (MVC)}

Es uno de los patrones aquitectónicos \cite{web:mvc} que más se utilizan en el desarrollo de aplicaciones, al tener tres capas diferenciadas cuyo objetivo es separar la vista del modelo de los datos subyacentes.

El patrón divide la aplicación en las siguientes capas:

\begin{itemize}
	\item Modelo: Se corresponde con el acceso a los datos, se encarga de almacenar y gestionar los que utiliza la aplicación.
	\item Vista: Capa de presentación, es la parte que ve el usuario, es la visualización del contenido de los datos del modelo.
	\item Controlador: Es la capa encargada de relacionar las dos capas anteriores. 
	\begin{itemize}
		\item Envía los datos del modelo a las vistas para que accedan los usuarios.
		\item Envía los eventos de los usuarios a los modelos, para actualizar los datos o traer otros nuevos.
	\end{itemize}
\end{itemize}

\imagenB{diagramaMVC}{Patrón MVC.}{.5}

\subsection{Repositorio}

Para gestionar parte de los modelos de datos, se ha implementado un patrón de repositorio, que se encarga de implementar la lógica de la gestión de las clases del modelo implicadas con la base de datos \textit{mysql}.

El resto de datos que utiliza nuestra aplicación provienen de archivos \textit{JSON} que contienen los cuestionarios.

\subsection{Diseño de paquetes}

En la elección del diseño de paquetes se ha tenido en cuenta la funcionalidad de cada clase, para agruparlas por características.

La aplicación consta de 5 paquetes diferenciados:

\begin{itemize}
	\item \textit{Config}: Se corresponde con la configuración de la aplicación.
	\item \textit{Controller}: Se corresponde con la capa de \textit{controlador} del modelo MVC.
	\item \textit{Model}: Se corresponde con la capa de \textit{modelo} del modelo MVC.
	\item \textit{Repository}: Contiene las clases necesarias para conectar los modelos con la base de datos.
	\item \textit{Services}: Contiene diferentes utilidades utilizadas en las diferentes clases de la aplicación.
\end{itemize}

\subsection{Arquitectura general}

La arquitectura general de la aplicación consta de dos contenedores independientes, uno con la base de datos y otro con la aplicación web.

El contenedor de la aplicación web accede al de la base de datos para comprobar si los usuarios que inician sesión están registrados y los datos son correctos.

A parte, el contenedor de la aplicación permite la conexión a la aplicación a través del puerto configurado.

La estructura quedaría de la siguiente forma:
\imagenB{diagramaAquitectura}{Diagrama de arquitectura.}{.4}



\apendice{Documentación técnica de programación}

\section{Introducción}

En el anexo actual se va a describir toda la documentación técnica, así como la estructura de directorios que se ha seguido para almacenar el proyecto completo.
Una de las partes más importantes es la sección de \textit{Manual del Programador} ya que en ella se describen los pasos necesarios para instalar el entorno de desarrollo. 

\section{Estructura de directorios}

El repositorio del proyecto está formado por los siguientes directorios:

\begin{itemize}
	\item \textbf{/}: Contiene el archivo README, la copia de la licencia y el comprimido riskteal.zip, con la versión corta para la instalación de la aplicación.
	\item \textbf{/docs}: Contiene la memoria y anexos en formato .pdf y .tex y la bibliografía en formato .bib.
	\item \textbf{/docs/tex}: Contiene lo diferentes archivos.tex que se han utilizado en la generación de los anexos y la memoria.
	\item \textbf{/docs/img}: Contiene las imágenes utilizadas en la construcción de los anexos y la memoria.
	\item \textbf{/docs/javadoc}: Contiene la documentación de la aplicación en formato javadoc.
	\item \textbf{/riskreal}: Contiene el docker-compose.yml que se utiliza para la instalación de la aplicación así como las carpetas \textbf{dbdata} y \textbf{appFiles} necesarias para el proceso.
	\item \textbf{/java}: Contiene las carpetas necesarias para el funcionamiento del entorno de desarrollo de la aplicación, así como todo el código completo del proyecto.
	\item \textbf{/java/csvFiles}: Carpeta donde se guarda el score.csv con los resultados de los cuestionarios.
	\item \textbf{/java/images}: Carpeta donde se guardan las imágenes de los cuestionarios.
	\item \textbf{/java/jsonFiles}: Carpeta que contiene los cuestionarios pendientes de generar.
	\item \textbf{/java/jsonFiles/lang}: Carpeta que contiene los cuestionarios ya generados en la aplicación.
	\item \textbf{/java/mysql}: Carpeta que contiene la carpeta \textbf{dbdata}, que se utiliza como ubicación de persistencia de la base de datos del entorno de desarollo.
	\item \textbf{/java/riskrealApp}: Carpeta que contiene el proyecto de maven y Spring Boot completo de la aplicación.
	\item \textbf{/java/riskrealApp/src/main}: Contiene todas las clases principales de java necesarias para que funcione la aplicación.
	\item \textbf{/java/riskrealApp/src/resources}: Contiene los archivos de configuración de la aplicación y el fichero quiz\_schema.json, con la estructura que deben seguir todos los json que se vayan a cargar.
	\item \textbf{/java/riskrealApp/src/resources/i18n}: Contiene los archivos de internacionalización de la aplicación.
	\item \textbf{/java/riskrealApp/src/resources/templates}:  Contiene las plantillas html necesarias para que funcione cada menú de la aplicación.
	\item \textbf{/java/riskrealApp/src/resources/static}: Contiene los ficheros de personalización para que se muestren las páginas web en un formato correcto. 
\end{itemize}

\section{Manual del programador}

El manual del programador es la guía \textit{paso a paso} de las diferentes herramientas, configuraciones y peculiaridades que se deben seguir para poder montar el entorno de desarrollo, trabajar con el código fuente o recompilar una nueva versión de la aplicación.

Las imágenes que se pueden ver en los siguientes pasos han sido tomadas desde un portátil con el sistema operativo Windows 11 y con las versiones aquí indicadas de las diferentes aplicaciones, por lo tanto, en distintas versiones el contenido o la ubicación de algunas opciones pueden estar en diferentes apartados.

\subsection{Herramientas necesarias}

Para poder desplegar el entorno de desarrollo se necesita instalar las siguientes aplicaciones:

\begin{itemize}
	\item Java SE JDK 17.
	\item Docker.
	\item Git.
	\item Spring Tools 4 for Eclipse.
\end{itemize}

A contunación, se explica cómo instalar y configurar cada uno de ellos de forma correcta.

\subsection{Java SE JDK 17}

Es el lenguaje de programación utilizado para la realización de la aplicación web.
Se debe acceder al siguiente enlace \cite{web:JavaJDK}, seleccionar el sistema operativo y arquitectura correspondiente y seguir el asistente de instalación.

\subsection{Docker}

En el momento actual en el que se encuentra la informática y el interés por la optimización de recursos, Docker se ha posicionado como la plataforma principal de \textit{microservicios}.

Docker Desktop se puede descargar desde el enlace \cite{web:dockerDesktop}. Tiene versión para windows, linux y mac, por lo que hay que elegir el sistema operativo correcto y seguir los pasos del asistente de instalación.

\textbf{La instalación exige un reinicio del equipo}, que se recomienda hacer en este punto para poder seguir con el resto del manual de forma secuencial.

Una vez instalado, la aplicación sería similar a la siguiente:
\imagen{dockerDesktopMain}{Pantalla principal de Docker Desktop.}

Podemos cerrar la aplicación sin problema porque se queda en ejecución en segundo plano, pero para confirmarlo, podemos abrir \textit{el símbolo del sistema - cmd} o una consola de powershell y ejecutar el comando \textbf{docker version}, donde deberíamos obtener una imagen similar a la siguiente:
\imagen{dockerVersion}{Versión de docker en ejecución desde CMD de Windows.}

Lo importante de la imagen anterior es que aparezca la información de \textbf{Server: Docker Desktop}, si no apareciese, abrir la aplicación de Docker Desktop y comprobar que ahora si tenemos la información del servidor.

\subsection{Git}

Git es la herramienta que nos permite recuperar todo el proyecto desde el repositorio web \cite{github:proyect}. 
Nos descargamos Git desde el siguiente enlace \cite{web:git}.
Seleccionamos la versión que se corresponda con nuestro sistema operativo y seguimos los pasos del asistente de instalación.

Una vez instalado, vamos a utilizar la herramienta \textbf{Git Bash}
\imagen{gitBash}{Aplicación Git Bash.}

Navegamos por los directorios hasta el directorio en el que queramos replicar el proyecto y a continuación, escribimos \textit{git clone https://github.com/Daniel-Fernandez-UBU/riskRealUBU}
\imagen{gitClone}{Comando Git Clone.}

\subsection{Spring Tools 4 for Eclipse}

\textit{Spring Tools 4 for Eclipse} es la herramienta desde la que se ha construido el proyecto de Spring al completo. 

Se puede obtener desde el siguiente enlace \cite{web:springtoolssuite}
\imagen{springTools4forEclipse}{Web de descarga de Spring Tools 4.}

En el caso de Windows, el fichero descargado es un \textit{.jar}. 

\begin{enumerate}
	\item Hacemos doble click sobre el fichero .jar descargado, para que se inicie la extracción automática del contenido.
	\begin{itemize}
		\item Si lo anterior no funciona, se tiene que realizar \textit{java -jar ``nombre-archivo-jar.jar''} para que se inicie.
	\end{itemize}
	\item En la carpeta en la que esté el fichero .jar se generará una nueva carpeta con nombre \textbf{sts-4.22.1.RELEASE}, o similar, dependiendo de la versión de la aplicación en el momento de la descarga.
	\item Movemos esa carpeta a la ubicación que nos interese, pues el programa de Spring Tools es \textit{portable}.
	\item Accedemos a la carpeta y ejecutamos \textbf{SpringToolSuite4.exe}
\end{enumerate}

Tras unos instantes en los que aparece la siguiente ventana:
\imagen{springToolsStart}{Spring Tools Suite 4 inicio.}

Nos pedirá que indiquemos la ubicación de nuestro espacio de trabajo, donde se guardarán por defecto todos los proyectos que creemos desde la aplicación.
\imagen{springToolsWork}{Spring Tools Suite 4 espacio de trabajo.}

Tras indicar nuestra ubicación preferida, pulsamos en \textit{Launch}.

Tras unos instantes en los que carga la aplicación, se nos abre la aplicación:

\imagen{springToolsIDE}{Spring Tools Suite 4 - IDE.}

Antes de importar el proyecto, vamos a instalar algunos plugins que pueden resultar de utilidad.

El proceso para la instalación de los plugins es el mismo, solo cambia el plugin a buscar.

Pasos a seguir para la instalación de cualquier plugin:
\begin{enumerate}
	\item Pulsamos en \textit{Help} --> \textit{Eclipse Marketplace\dots}
	\item Introducimos el nombre del plugin que queremos buscar.
	\item Pulsamos en \textit{Go} o le damos al \textit{Intro} en el teclado.
	\item Pulsamos en \textit{Install}.
	\item Seguimos los pasos, aceptando las diferentes ventanas de confirmación que nos aparecen.
	\item Tras terminarse la instalación del plugin, nos pide reiniciar Spring Tools Suite 4. \textbf{Consejo: Instalar todos los plugins y al final reiniciar la aplicación}
\end{enumerate}

\subsubsection{Thymeleaf}

Siguiendo los pasos anteriores, obtendríamos la siguiente ventana, donde ya solo quedaría darle a \textit{Install} y seguir los pasos del asistente de instalación.
\imagen{pluginThymeleaf}{Instalación del plugin de Thymeleaf en Spring Tools Suite.}

\subsubsection{Eclipse Web Developer}

Instalamos la versión \textit{Eclipse Enterprise Java and Web Developer Tools 3.33} que incluye los editores de html, json y css entre otros.
\imagen{pluginEnterpriseJava}{Instalación del plugin de Eclipse Enterprise Java and Web en Spring Tools Suite.}

Seleccionamos las características que vamos a instalar:
\imagen{pluginEnterpriseJavaFeatures}{Selección de características de Eclipse Enterprise Java and Web en Spring Tools Suite.}


\subsubsection{Eclipse Docker Tooling}

Este plugin viene instalado por defecto en la última versión disponible de Spring Tools 4 for Eclipse, pero si no apareciese en la sección de \textit{Installed} en \textit{Eclipse Marketplace}, también sería necesario instalarlo.

\imagen{pluginDocker}{Plugin de Docker en Spring Tools Suite.}


\subsubsection{Importación de nuestro proyecto}

Una vez hayamos finalizado con la instalación de los plugins anteriores, procederemos a importar nuestro proyecto.

En primer lugar seleccionamos \textbf{File --> Import\dots}

\imagen{fileImport}{Importar proyecto en Spring Tool Suite.}

En la nueva ventana que se muestra, navegamos hasta \textbf{Maven --> Existing Maven Projects --> Next}.

\imagen{importMaven}{Importación de proyecto Maven en Spring Tool Suite.}

A continuación, nos pedirá que busquemos la ruta donde tenemos el proyecto que queremos importar. \textbf{Tenemos que ir a la ruta sobre la que hayamos ejecutado el \textit{git clone} del paso previo}.

Accedemos a la ruta donde hemos hecho el clon del repositorio de Git, y dentro de la carpeta del proyecto, \textit{riskRealUBU}, accedemos a \textit{java --> riskRealApp}.

Si hemos seleccionado la ruta adecuada, deberá quedar de esta forma la importación, habiendo reconocido el fichero \textbf{pom.xml} de nuestro proyecto y a la espera de que pulsemos en \textit{Finish}.

\imagen{browseProject}{Ruta del proyecto indicada en Spring Tool Suite.}

Tras unos instantes, ya nos aparecería el proyecto correctamente cargado en nuestra aplicación.

\imagen{proyectImported}{Proyecto importado en Spring Tools Suite.}

\subsubsection{Iniciando la aplicación}

Tras haber completado correctamente todos los pasos anteriores, y teniendo \textit{Docker Server} funcionando, desde Spring Tools Suite podemos ejecutar la aplicación.

\textbf{Los puertos por defecto que usa son: 3306 para mysql y 8088 para la aplicación web, sería recomendable no tener ninguno de los 2 en uso, para el correcto despliegue de la aplicación sin necesidad de tener que cambiar opciones de configuración en \textit{application.properties}.}

En la parte inferior izquierda de la interfaz de Spring Tool, tendremos \textit{Boot Dashboard} y si desplegamos \textit{local} nos aparecerá nuestra aplicación \textbf{riskrealApp}

Seleccionamos nuestra aplicación y pulsamos en el icono que tenemos justo encima con un \textit{cuadrado rojo y triángulo verde} para iniciarla.

\imagen{startApp}{Iniciar la aplicación en Spring Tools Suite.}

Tras iniciar la aplicación, si es la primera vez tardará un poco más porque tiene que descargarse de docker la imagen de MySQL, debería quedarse en un estado similar a este:

\imagen{appStarted}{Aplicación iniciada en Spring Tools Suite.}


Ya tenemos nuestro entorno de desarrollo configurado, con nuestra aplicación iniciada y accesible para poder probarla desde \url{http://localhost:8088}.

\section{Compilación, Generación de una nueva imagen, y Docker Compose}

\subsection{Compilación}

Para compilar el proyecto y generar una nueva imágen de docker de nuestra aplicación, tenemos que realizar los siguientes pasos:

\begin{enumerate}
	\item Ejecutar desde Spring Tools Suite 4: Botón derecho del ratón en el nombre del proyecto --> Run as\dots --> Maven build\dots --> Escribimos \textbf{spring-boot:build-image}
	\item Una vez termine ya tenemos la imagen cargada en docker --> Se puede comprobar con \textbf{docker images} desde el símbolo del sistema o en la aplicación de Docker Desktop.
\end{enumerate}

\subsection{Generación de una nueva imagen}

Tras lo anterior, podemos subir la imagen a nuestro repositorio de github de forma muy sencilla siguiendo estos pasos:

\begin{enumerate}
	\item Renombramos la imagen de docker siguiendo el patrón: \textbf{ghcr.io/<username de github>/<repositorio>/<nombre de imagen>:<version>} (la versión es opcional).
 	\item Por ejemplo: \textbf{docker image tag riskrealapp:latest ghcr.io/daniel-fernandez-ubu/riskrealubu/riskrealapp}
 	\item Crear un \textbf{Access Token (classic)} en nuestro perfíl de GitHub --> Settings --> Developer Settings (Está abajo del todo)
 	\item Iniciar sesión en la terminal (o línea de comandos) desde donde vayamos a hacer el push con el comando:
 	\item \textbf{echo <token> | docker login ghcr.io -u USERNAME --password-stdin}
 	\item Hacer pull de nuestra imagen al repositorio:
 	\item \textbf{docker push ghcr.io/daniel-fernandez-ubu/riskrealubu/riskrealapp}. Al no poner la versión, utiliza \textit{latest} por defecto.
	\item Si hemos vinculado el repositorio de \textbf{package} con nuestro repositorio, ya al acceder vemos que hay una nueva imagen.
\end{enumerate}

\subsection{Fichero docker-compose.yml}\label{subsec:dockerComposeFile}

Es el fichero más importante de cara a lanzar una nueva versión de la aplicación, por ese motivo se ha decidido hacer una pequeña introduccion del contenido del fichero, para entender el motivo de cada línea y para que cada programador pueda configurarlo a su gusto.

El contenido del fichero es el siguiente:
\imagen{dockerComposeYML}{Contenido del fichero docker-compose.yml}

A continuación, explicamos cada una de sus líneas:
\begin{itemize}
	\item services: --> Etiqueta que se utiliza para definir los servicios o dockers que se van a crear.
	\begin{itemize}
		\item db: --> Servicio para el docker que contiene la base de datos.
		\begin{itemize}
			\item image: mysql --> Imagen de la base de datos que se va a utilizar.
			\item container\_name: riskreal-db-1 --> Nombre del docker que se va a crear.
			\item restart: always --> Reinicia el docker si se genera un error.
			\item environment: --> Definimos varias variables de entorno de la imagen que vamos a lanzar.
			\begin{itemize}
				\item MYSQL\_DATABASE: appRiskRealUbu --> Nombre de la base de datos que se va a crear.
				\item MYSQL\_ROOT\_PASSWORD: rootroot --> Contraseña de root para acceder a la base de datos.
			\end{itemize}
			\item expose: --> Puertos internos que se van a abrir, para que el resto de contenedores de la misma red puedan acceder.
			\begin{itemize}
				\item - '3306' --> Se habilita el puerto 3306 para acceder a este docker por la red interna.
			\end{itemize}
			\item ports: --> Puertos que se publican del docker para acceder a él desde fuera de la red interna, es decir, desde el equipo anfitrión.
			\begin{itemize}
				\item '30306:3306' --> Permite el acceso a la base de datos desde el puerto 30306.
			\end{itemize}
			\item volumes: --> Se definen los directorios que queremos que sean persistentes, es decir, que tras el reinicio o eliminación del docker, se mantengan.
			\begin{itemize}
				\item ./dbdata:/var/lib/mysql --> Se utiliza la carpeta local dbdata para almacenar la base de datos completa. Si dbdata tuviese una base de datos MySQL, el docker la utilizaría al iniciar.
			\end{itemize}
			\item networks: --> Define la red interna del docker.
			\begin{itemize}
				\item - riskreal-network --> Nombre de la red interna.
			\end{itemize}
			\item healtcheck: --> Comprobación adicional para saber cuando ha terminado de levantar la base de datos en el docker.
			\begin{itemize}
				\item test: --> Se lanza un ping con mysqladmin a localhost para saber cuando ha iniciado la base de datos.
				\item interval: 30s --> Define el intervalo de tiempo en el que se comprueba el estado.
				\item timeout: 10s --> Indica el intervalo de tiempo que espera antes de indicar que no hay respuesta del servidor.
				\item retries: 5 --> Número de veces que lo reintenta antes de indicar que el docker no ha iniciado correctamente.
			\end{itemize}
		\end{itemize}
		\item webapp: --> Servicio para el docker que contiene la aplicación web.
		\begin{itemize}
			\item image: ghcr.io/daniel-fernandez-ubu/riskrealubu/riskrealapp:v1 --> Imagen de docker que se usará. Imagen propia de la Spring Boot App.
			\item container\_name: riskreal-webapp-1 --> Nombre del docker que se va a crear.
			\item environment: --> Definimos varias variables de entorno de la imagen que vamos a lanzar.
				\begin{itemize}
				 	\item JSON\_QUIZ\_FILE\_PATH: /opt/jsonFiles --> Ruta en la que se guardan los json con los cuestionarios que se quieren cargar en la aplicación.
					\item JSON\_QUIZ\_FILE\_PATH\_LANG: /opt/jsonFiles/lang --> Ruta en la que se guardarán los cuestionarios ya cargados en la aplicación.
					\item CSV\_SCORE\_PATH: /opt/csvFiles/scores.csv --> 
- Ruta y nombre del fichero de resultados que se rellenará conforme se vayan realizando cuestionarios.
					 \item SPRING\_DATASOURCE\_URL: \\ jdbc:mysql://db:3306/appRiskRealUbu --> Dirección de la base de datos. 
					 \item SPRING\_DATASOURCE\_PASSWORD: rootroot --> Contraseña de root para acceder a la base de datos.
					 \item SPRING\_DATASOURCE\_DRIVER\_CLASS\_NAME: \\com.mysql.cj.jdbc.Driver --> Driver de conexión a la base de datos.
					 \item SERVER\_PORT: 8088 --> Puerto en el que se publicará la aplicación.
				 	 \item ADMIN\_USERNAME: admin@riskreal.com --> Usuario administrador por defecto.
				 	 \item ADMIN\_PASSWORD: admin --> Contraseña por defecto del usuario administrador.
  				 	 \item TEST\_QUIZ: test\_quiz\_ --> Nombre por defecto de los archivo json de test para usuarios anónimos.
				\end{itemize}
			\item ports: --> Puertos que se publican del docker para acceder a él desde fuera de la red interna, es decir, desde el equipo anfitrión.
			\begin{itemize}
				\item '8088:8088' --> Redirecciona el puerto 8088 del servidor anfitrión al puerto 8088 interno del docker, en el que se ha publicado la aplicación.
			\end{itemize}
			\item volumes: --> Se definen los directorios que queremos que sean persistentes, es decir, que tras el reinicio o eliminación del docker, se mantengan.
			\begin{itemize}
				\item ./appFiles:/opt --> Se utiliza para indicar la ruta de los cuestionarios que queremos cargar, los ya almacenados y el fichero de resultados.
				\item ./appFiles/images:/workspace/BOOT-INF/classes/static/images --> \\Se utiliza para indicar la ruta donde tenemos las imágenes de los cuestionarios.
			\end{itemize}
			\item networks: --> Define la red interna del docker.
			\begin{itemize}
				\item - riskreal-network --> Nombre de la red interna.
			\end{itemize}
			\item depends\_on: --> Se definen los servicios de los que depende el inicio de este docker.
			\begin{itemize}
				\item db: --> Servicio del que depende.
				\item condition: service\_healthy --> Necesita que se haya iniciado de forma correcta la base de datos para iniciar el docker.
			\end{itemize}
		\end{itemize}
	\end{itemize}
	\item neworks: --> Sección en la que se definen las redes.
	\begin{itemize}
		\item riskreal-network: --> Nombre de la red definida.
	\end{itemize}
\end{itemize}

\section{Pruebas del sistema}

Al tratarse de una aplicación web, no ha requerido el uso y/o desarrollo de una batería de pruebas.

La mayoría de las funcionalidades se basan en el registro o modificación de datos de usuarios, donde los propios campos de los html tienen las restricciones necesarias para que se cumpla con lo esperado.

Las únicas funcionalidades que podrían requerir de pruebas serían:

\begin{itemize}
	\item La carga de un archivo json desde el sistema de archivos a la aplicación.
	\begin{itemize}
		\item Si no se selcciona un archivo json la aplicación te avisa y no deja cargarlo.
		\item Si se selecciona un archivo json pero su estructura es incorrecta, no lo carga y avisa.
	\end{itemize}
	\item El envío de un correo electrónico de recuperación de contraseñas.
	\begin{itemize}
		\item Si el correo electrónico introducido está en la base de datos, se envía el email.
		\item Si no lo está, se informa indicando que no es correcto.
		\item Si ha sucedido algún error que no permita enviar el correo, se informa indicando de esta situación: Solo podría darse si hubiese problemas con el servidor de correo o la cuenta utilizada para el envío.
	\end{itemize}
\end{itemize}

\apendice{Documentación de usuario}

\section{Introducción}

Este manual detalla todos los pasos necesarios para poder instalar en cualquier plataforma la aplicación. Se definirán los requisitos previos que tiene que tener los usuarios así como un manual de uso de la propia página web.

\section{Requisitos de usuarios}

La aplicación utiliza la tecnología de docker para funcionar, con todas sus ventajas y un pequeño inconveniente:

\begin{itemize}
	\item \textbf{Necesitas tener docker instalado en el equipo en el que vayas a desplegar la aplicación.}
\end{itemize}

Los requisitos indispensables para poder utilizar la aplicación son los siguientes:

\begin{enumerate}
	\item Instalar la plataforma de docker \cite{web:dockerDesktop} y docker compose \cite{web:dockercompose}.
	\begin{itemize}
		\item Windows: Docker Desktop.
		\item Linux: Instalar Docker y Docker Compose, ya sea con Docker Desktop o desde el repositorio de la versión de linux correspondiente.
		\item Mac: Docker Desktop.
	\end{itemize}
	\item Disponer de un navegador web.
	\item Disponer de un programa para descomprimir ficheros \textit{.zip} o \textit{.tar.gz}.
\end{enumerate}

\section{Instalación}

\subsection{Descarga del código}

La descarga del código se realiza desde esta página: \url{https://github.com/Daniel-Fernandez-UBU/riskRealUBU/releases}

\imagen{instalacionReleases}{Página de selección de release}

Lo recomendable es siempre utilizar la última versión, ya que tendrá el mayor número de funcionalidades y el menor número de errores.

De la imagen anterior se puede ver que tenemos la posibilidad de descargar tres posibles ficheros:

\begin{itemize}
	\item Source code: Comprimido en dos formatos distintintos: Contiene el código completo del proyecto.
	\item Install.zip: Contiene la carpeta de instalación, la única que vamos a necesitar para poner a funcionar nuestra aplicación.
\end{itemize}

\textit{Nota: Si has descargado Source code, tras descomprimirlo y acceder a la carpeta del proyecto, tendrás una carpeta install, que es donde está la parte que importa para la instalación de la aplicación.}

\subsection{Contenido de la carpeta install}

La estructura de la carpeta es la siguiente:

\imagen{installTree}{Contenido y estructura de la carpeta install.}

\subsubsection{Fichero docker-compose.yml}

Este es único fichero realmente necesario, ya que con el simple comando \textbf{docker compose up} en el mismo directorio en el que se encuentra este fichero, se levanta la aplicación.

La explicación del contenido se ha realizado en \ref{subsec:dockerComposeFile}.

\subsection{Instalación de la aplicación}

Desde la línea de comandos, navegamos hasta acceder a la carpeta install y estar al mismo nivel que \textit{docker-compose.yml, dbdata y appFiles}.

\subsubsection{Instalación supervisada}

Escribimos \textbf{docker compose up} y empezamos a ver cómo se lanzan las aplicaciones paso a paso. Si es la primera vez, tiene que hacer una descarga de las imágenes al repositorio de imágenes de docker, por lo que es normal que lleve algo más de tiempo.
\imagen{dockerComposeUP}{Comando para la instalación.}

En la siguiente imagen se ve el resultado de una ejecución correcta:
\imagen{dockerComposeRunning}{Aplicación funcionando.}

Desde otra consola de símbolo del sistema o desde Docker Desktop, podemos ver nuestros docker funcionando:
\imagen{dockerWorking}{Dockers en ejecución.}

Para parar la ejecución y eliminar los docker, pulsamos \textit{Ctrl + C} y escribimos \textit{docker compose down} en la ventana del símbolo del sistema en la que los teníamos en ejecución.
\imagen{dockerStop}{Parada de los docker lanzados.}

\subsubsection{Instalación desatendida}

El proceso es el mismo que en el paso anterior, con la diferencia de que esta vez no vamos a ver cómo se van lanzando los contenedores, pero con la ventaja de que vamos a poder cerrar la ventana del símbolo del sistema y van a seguir en ejecución.

Escribimos \textbf{docker compose up -d} y comprobamos con \textit{docker ps} cuando se han levantado nuestros dockers:
\imagen{dockerComposeUPD}{Comando para la instalación desatendida.}

Si queremos parar la aplicación podemos lanzar el comando \textit{docker compose down} explicado en la sección anterior, pero tenemos que lanzarlo estando a la altura de nuestro fichero \textit{docker-compose.yml} para que funcione.

\section{Manual del usuario}

Este manual va a servir como guía para poder navegar por la aplicación de forma sencilla.
Nota: \textit{Las imágenes de los cuestionarios que se muestran en este manual son imágenes de prueba que se han precargado en cuestionarios de prueba, aunque siempre aparezca la misma.}

En primer lugar accedemos a la aplicación, actualmente la url por defecto de acceso es \url{localhost:8088}.

\imagen{appInicio}{Pantalla de inicio de la aplicación web.}

En la barra superior tenemos varias opciones disponibles, la de cambio de idioma:

\imagen{appIdioma}{Selección de idioma de la aplicación web.}

Y la de Inicio de Sesión o Registro:

\imagen{appInvitado}{Opciones de usuario de la aplicación web.}

También se puede realizar el cuestionario disponible para usuarios anónimos, pulsando en el botón que tenemos en el centro de la imagen del cuestionario.

Si seleccionas Registro, te saldrá un formulario con diversos campos que tienes que completar, algunos de ellos solo permiten valores de un listado:

\imagen{appRegistro}{Registro en la aplicación web.}

Una vez te registres, te envía a la página de inicio de sesión para que introduzcas tus datos y accedas a la aplicación.

Si en la página inicial hubieses seleccionado iniciar sesión:

\imagen{appInicioSesion}{Pantalla de inicio de sesión de la aplicación web.}

Te encontrarás con un formulario para incluir tu nombre de usuario (correo electrónico) y y tu contraseña. 

Tras hacerlo te devuelve a la página inicial, desde la que podrás iniciar un cuestionario, cambiar el idioma, cerrar sesión o ver tu perfíl.

Al ser un usuario autenticado, te saldrán estas dos últimas opciones que te comentaba al pulsar en tu nombre, en la barra de navegación de la aplicación:

\imagen{appUserEng}{Pantalla de inicio de la aplicación web en inglés.}

Se ha cambiado el idioma a inglés solo para mostrar la funcionalidad.

Si accedes a tu perfíl, podrás ver todos los datos que acabas de incluir, editar los que consideres e incluso cambiar tu contraseña:

\imagen{appPerfil}{Pantalla de perfíl de usuario.}

Si decides iniciar uno de los cuestionarios disponibles desde la página de inicio, te saldrá una pantalla similar a esta:

\imagen{appDesCuestionario}{Pantalla de descripción del cuestionario.}


Donde verás la imágen princiapl del cuestionario junto con una descripción del mismo. Finalmente podrás darle a iniciar cuando consideres.

Tras iniciar el cuestionario, tendrá un formato similar al de esta imagen:

\imagen{appIniCuest}{Pantalla de inicio de la aplicación web.}

Con las opciones a la izquierda, la imagen relacionada con la pregunta a la derecha y debajo 3 botones: Anterior, Siguiente o Enviar resultados.

Los botones, según avances, en las preguntas, se irán habilitando, para que puedas retroceder sin problema en cualquier momento, lo único indispensables es que si quieres cambiar de pregunta debes seleccionar alguna respuesta antes de poder hacerlo. Lo mismo sucede al enviar resultados, por lo que todas las preguntas deben haber sido seleccionadas.

Al llegar a la última pregunta y pulsar en ``Enviar Resultados'', te saldrá una nueva pantalla, con el resultado que has obtenido en el test:

\imagen{appResultado}{Pantalla de resultado del cuestionario.}

Si has quieres hacer el cuestionario siendo un invitado, la única diferencia es que tras pulsar en el cuestionario disponible en la página de inicio, antes del proceso anterior, te sale una pantalla donde se te piden algunos datos adicionales:

\imagen{appDataGuest}{Pantalla de datos adicionales para los usuarios anónimos.}

Por último, si has iniciado sesión con un usuario administrador, tendrás acceso a más opciones de las permitidas hasta ahora:

\imagen{appAdmin}{Opciones del usuario administrador.}

Vamos a comentar cada una de ellas:

Al seleccionar en la opción de ``Cargar Cuestionario'' accedemos a la siguiente pantalla, desde la que podemos o bien cargar un nuevo json desde el disco duro de nuestro equipo:

\imagen{appCargarJSON}{Pantalla de carga de archivos json.}

O bien generar un cuestionario utilizando un json previamente cargado:

\imagen{appGenerarCuestionario}{Pantalla de generación de cuestionarios.}

La opción de ``Descargar resultados'' la tenemos disponible desde cualquier parte de la web, excepto mientras se está realizando un cuestionario, y tras pulsar en ella, se nos descarga de forma automática un archivo \textit{score.csv} con los datos de todos los cuestionarios que se han realizado.

\imagen{appDescargaResultados}{Pantalla tras pulsar en descargar resultados.}

El usuario administrador, desde la opción de Gestión de Roles, puede modificar los permisos sobre la aplicación de otros usuarios, y convertirlos también en administrador:

\imagen{appRoles}{Pantalla de gestión de roles de la aplicación web.}

El rol del usuario administrador no es modificable en ningún caso.

Si al intentar iniciar sesión, te equivocas en la contraseña, te saldrá una nueva ventana donde puedes probar de nuevo o restarurarla:
\imagen{appLoginError}{Pantalla inicio de sesión erróneo.}


Si pulsas en \textit{Recuperar contraseña} se abre una nueva ventana donde se te pide el el correo electrónico.
\imagen{appPassRecovery}{Pantalla inicio de recuperación de contraseña.}

Si el correo introducido es correcto y el usuario está registrado en la base de datos de la aplicación, pasados unos segundos le llegará un correo a su buzón y la aplicación te devolverá a la página de inicio.

El correo de recuperación que se recibe es similar a este:
\imagen{appEmailReceived}{Correo con la nueva contraseña de un solo uso.}

Tras iniciar sesión con la contraseña que has recibido, te pedirá que la cambies:
\imagen{appCambioPass}{Pantalla cambio de contraseña.}

Después de cambiarla, volverás a la página de inicio con la sesión iniciada.

\clearpage
Por último, si deseas salir de aplicación, al pulsar en ``Cerrar Sesión'' te pedirá confirmación antes de hacerse efectivo:

\imagen{appCerrarSesion}{Pantalla de cierre de sesión.}






\apendice{Anexo de sostenibilización curricular}

\section{Introducción}
Este anexo incluirá una reflexión personal del alumnado sobre los aspectos de la sostenibilidad que se abordan en el trabajo.
Se pueden incluir tantas subsecciones como sean necesarias con la intención de explicar las competencias de sostenibilidad adquiridas durante el alumnado y aplicadas al Trabajo de Fin de Grado.

Más información en el documento de la CRUE \url{https://www.crue.org/wp-content/uploads/2020/02/Directrices_Sosteniblidad_Crue2012.pdf}.

Este anexo tendrá una extensión comprendida entre 600 y 800 palabras.



\bibliographystyle{plain}
\bibliography{bibliografiaAnexos}

\end{document}
